% Save as: module10-problem-set.tex
% ============================================================================
% Problem Set 10: AI-Econ Nexus
% Microeconomic Theory for Experimentalists & Applied Microeconomists
% Built from templates/problem_set_template.tex. Compile with pdflatex.
% ============================================================================
\documentclass[11pt]{article}
\usepackage{amsmath, amssymb, amsthm}
\usepackage[margin=1in]{geometry}
\usepackage{enumitem}
\usepackage{booktabs}

\newtheorem{question}{Question}

\title{Problem Set 10: AI-Econ Nexus\\
\large Microeconomic Theory for Experimentalists \& Applied Microeconomists}
\author{Due: \textbf{[date --- before finals week]}}
\date{}

\begin{document}
\maketitle

\noindent\emph{Ground rules: collaboration on ideas is encouraged; write-ups
are individual. You may use AI tools for parts flagged with
$\langle$AI$\rangle$ and must attach the transcript. Show all calculations.}

\medskip

\noindent\textbf{Weights:} Part A --- 30 points. Part B --- 45 points.
Part C --- 25 points.

\bigskip

% ---------------------------------------------------------------------------
\section*{Part A: Theory --- Guided Calculations (30 points)}

\begin{question}[What is a replication worth? Bayes on silicon evidence
--- 15 points]
You hold a prior of $P(F) = 0.3$ that your agent population is a
\emph{faithful simulator} of the target human population (versus
$\neg F$: it produces answers by recitation and surface patterns
without a faithful behavioral process).

\begin{enumerate}[label=(\alph*)]
  \item \textbf{(5 pts)} You run a famous experiment and the agents
        reproduce the published result. Because the result saturates
        the training data, recitation predicts this success almost as
        well as fidelity: $P(\text{reproduce} \mid F) = 0.9$ and
        $P(\text{reproduce} \mid \neg F) = 0.85$. Compute the
        posterior $P(F \mid \text{reproduce})$ and state, in one
        sentence, what the small movement teaches.
  \item \textbf{(5 pts)} Starting from your posterior in (a), you now
        test a \emph{novel, preregistered, out-of-corpus} prediction
        (a design variation with no published answer):
        $P(\text{match} \mid F) = 0.7$,
        $P(\text{match} \mid \neg F) = 0.2$. Compute the posterior
        after a successful match. Why is this test so much more
        informative --- answer in terms of the likelihood ratio.
  \item \textbf{(3 pts)} Your labmate proposes boosting credibility by
        replicating THREE more famous results. Explain in two
        sentences why the posterior gain is small even after all
        three, and what single study design would move it more.
  \item \textbf{(2 pts)} One sentence: what does this calculation
        imply about how your capstone's simulation section should be
        structured?
\end{enumerate}
\end{question}

\begin{question}[Reading a stability report --- 15 points]
Your audit battery (Lab 10) runs a dictator-game pilot across two
models and three prompt paraphrases (same content, different wording),
at fixed temperature. Mean percent given:

\begin{center}
\begin{tabular}{@{}lccc@{}}
\toprule
 & Paraphrase 1 & Paraphrase 2 & Paraphrase 3 \\
\midrule
Model A & 27 & 29 & 41 \\
Model B & 33 & 34 & 36 \\
\bottomrule
\end{tabular}
\end{center}

For reference, human dictator-game means vary across comparable lab
replications by roughly $\pm 5$ percentage points around a common mean
--- treat spreads wider than about 10 points as clearly exceeding the
band.

\begin{enumerate}[label=(\alph*)]
  \item \textbf{(4 pts)} Compute each model's paraphrase range (the
        prompt-sensitivity index) and the cross-model gap at each
        paraphrase. Which numbers exceed the human cross-lab
        variability band?
  \item \textbf{(4 pts)} Classify the following claims as SURVIVES or
        DOES NOT SURVIVE this report, with one clause of
        justification each: (i) ``silicon giving is positive and
        substantial''; (ii) ``silicon mean giving is $\approx$30\%'';
        (iii) ``Model A exhibits human-like giving.''
  \item \textbf{(4 pts)} Diagnose Model A's Paraphrase 3: name two
        candidate explanations (one about the wording, one about the
        model) and the single follow-up run that would separate them.
  \item \textbf{(3 pts)} State the ``one model, one prompt'' fallacy
        in one sentence, and the minimal robustness protocol your
        capstone pilot must therefore include.
\end{enumerate}
\end{question}

% ---------------------------------------------------------------------------
\section*{Part B: Experimental Design (45 points)}

\begin{question}[The full sim-pilot workflow, pre-registered]
A research team plans a human field experiment: \emph{does paying gig
workers a fixed bonus upfront (versus the same bonus on completion)
change task quality?} The upfront bonus is unconditional --- paid at
the start and never clawed back. (Module 4's gift exchange and Module
5's reference dependence both make predictions; a clawback variant
would be a different design, pitting loss framing against gain
framing.) Before fielding it, they
will run an LLM simulation pilot. Design the complete pilot workflow.
Specify:
\begin{enumerate}[label=(\alph*)]
  \item \textbf{(6 pts)} The human design being piloted: task, arms,
        outcomes (quantity AND quality), and the two theories'
        competing predictions --- so the pilot has something specific
        to serve.
  \item \textbf{(12 pts)} The simulation pilot itself: agent design
        (personas, and what must NOT be smuggled into them), task
        rendering, treatments mirrored from the human design, and ---
        the heart --- the specific design parameters the pilot is
        allowed to tune (instruction wording, comprehension checks,
        bonus salience, quality-rubric clarity) versus the quantities
        it must not be used to estimate. Justify each assignment
        using the course's can/cannot list.
  \item \textbf{(9 pts)} The preregistration, written before any run:
        what gets committed (hypotheses, prompts, models,
        temperatures, N per cell, parsers, analysis code, the audit
        battery plan) and WHY cheap iteration makes this discipline
        more necessary, not less.
  \item \textbf{(9 pts)} The decision rule: state, in advance, the
        three result patterns --- one that justifies fielding the
        human study as designed, one that sends the team back to
        redesign (and what specifically gets redesigned), one that
        would halt the project --- and defend why a silicon result
        can legitimately trigger each.
  \item \textbf{(9 pts)} The steel-manned objection: write the
        STRONGEST version of ``this pilot tells you nothing --- the
        model has read your literature,'' then the response that
        concedes what is correct and states precisely what the pilot
        still delivers.
\end{enumerate}
\end{question}

% ---------------------------------------------------------------------------
\section*{Part C: Silicon Pilot Interpretation $\langle$AI$\rangle$
(25 points)}

\begin{question}[The committee member's objection]
At your capstone defense, a committee member says: ``I stopped reading
at the simulation section. These models are trained on the very
literature you study; your `pilot' is a book report with extra steps.
Nothing in it should influence whether the human study runs.''

\begin{enumerate}[label=(\alph*)]
  \item \textbf{(8 pts)} Steel-man the objection: restate it in its
        STRONGEST form, adding the two best supporting arguments the
        committee member did not think to make (draw on the course's
        divergence taxonomy --- contamination, alignment artifacts,
        missing channels, prompt fragility).
  \item \textbf{(9 pts)} The defense: concede explicitly what is
        right (which uses of your pilot the objection genuinely
        kills), then defend what survives --- the can-list uses ---
        showing for each WHY the objection's mechanism does not touch
        it. Your defense must cite at least one concrete result from
        your own audit battery (real if you ran it; specified
        precisely if not).
  \item \textbf{(8 pts)} Write the honest methods paragraph
        (4--6 sentences) for your capstone paper: what the simulation
        was used for, what it was not used for, the stability
        evidence, and where every effect-size and power assumption
        actually came from. This paragraph will be graded against the
        course's honesty standard: claiming a use from the
        cannot-list costs more than admitting a limitation.
\end{enumerate}
Attach your AI transcripts and, if you ran code, the results CSV.
\end{question}

\end{document}
